\clearpage\section{PERFIL PROFISSIONAL DO EGRESSO}
Segundo à Resolução CNE/CES nº 11, de 11 de março de 2002 que institui as Diretrizes Curriculares Nacionais do Curso de Graduação em Engenharia, especialmente quanto ao Art. 3°, a saber:

\begin{itemize}
	\item “O Curso de Graduação em Engenharia tem como perfil do formando egresso/profissional o engenheiro, com formação generalista, humanista, crítica e reflexiva, capacitado a absorver e desenvolver novas tecnologias, estimulando a sua atuação crítica e criativa na identificação e resolução de problemas, considerando seus aspectos políticos, econômicos, sociais, ambientais e culturais, com visão ética e humanística, em atendimento às demandas da sociedade”.
\end{itemize}

Neste escopo, o perfil do \NOMEPROFISSIONAL, formado através do trabalho interdisciplinar e do exercício prático dos conhecimentos adquiridos, resulta nas seguintes competências e habilidades gerais:

\begin{enumerate}
	\item Aplicar conhecimentos matemáticos, científicos, tecnológicos e instrumentais à engenharia;
	\item Projetar e conduzir experimentos e interpretar resultados;
	\item Conceber, projetar e analisar sistemas, produtos e serviços de engenharia;
	\item Participar de forma responsável, ativa, crítica e criativa nos processos de projetos de automação e
	controle;
	\item Identificar, formular e resolver problemas de engenharia desenvolvendo ou utilizando novas ferramentas e técnicas;
	\item Avaliar criticamente a operação e a manutenção de sistemas;
	\item Elaborar documentação técnica sobre equipamentos e sistemas de automação e controle;
	\item Pesquisar e desenvolver novas tecnologias na área de automação e controle;
	\item Atuar em equipes multidisciplinares na execução de projetos industriais de automação e controle;
	\item Planejar, supervisionar, elaborar e coordenar projetos e serviços de engenharia;
	\item Compreender e aplicar a ética e responsabilidade profissionais;
	\item Avaliar o impacto das atividades da engenharia no contexto social e ambiental, em especial as que são contempladas nos arranjos produtivos locais e regionais, tais como: têxtil, alimentícia, siderúrgica, química, manufatura e outras;
	\item Assumir a postura de permanente busca de atualização profissional, com vistas a capacitar-se e, assim, adaptar-se às novas tecnologias e demandas apresentadas pelo mundo do trabalho.
\end{enumerate}

O aluno egresso também poderá se dedicar ao desenvolvimento e gerência do próprio negócio, tornando-se um empreendedor. Para tanto, o engenheiro formado deverá ter sólida formação técnico científica e profissional. Esta formação o capacita a desenvolver novas tecnologias, estimulando a sua atuação crítica e criativa na identificação e resolução de problemas, considerando seus aspectos políticos, econômicos, sociais, ambientais e culturais, com visão ética e humanística, em atendimento às demandas da sociedade, desde os níveis local e regional até os níveis nacional e mundial. Alicerçado numa formação abrangente, ele estará capacitado para exercer ação integradora, podendo ser considerado como um Engenheiro de Sistemas orientado à concepção, implementação, uso e manutenção de sistemas automatizados. Sua formação diferencia-se, assim, daquela do engenheiro de processo (mecânico, químico, elétrico etc.). 



